\section*{الفصل \ref{ch:g12:equilibrium} --- التوازن الكيميائي: خارج التفاعل وثابت التوازن}

\begin{solution}{exo:g12:equilibrium:1}
(a) $Q = \dfrac{[\ce{NH3}]^2 (c^\circ)^2}{[\ce{N2}][\ce{H2}]^3}$؛\par
(b) $Q = \dfrac{[\ce{CH3COO-}][\ce{H3O+}]}{[\ce{CH3COOH}]\,c^\circ}$ (يُستبعد الماء،
وهو \omterm{def:g6:solutions-and-solubility:solute}{المذيب})؛\par
(c) $Q = [\ce{Ca^{2+}}][\ce{CO3^{2-}}]/(c^\circ)^2$ (يُستبعد الصلب)؛\par
(d) $Q = \dfrac{[\ce{FeSCN^{2+}}]\,c^\circ}{[\ce{Fe^{3+}}][\ce{SCN-}]}$.
\end{solution}

\begin{solution}{exo:g12:equilibrium:2}
$\tau = 0.020/0.050 = 0.40$: فهو غير تام.
\end{solution}

\begin{solution}{exo:g12:equilibrium:3}
$Q = 2 < K$: مباشر. $Q = 50 > K$: عكسي. $Q = K$: لا تغيّر.
\end{solution}

\begin{solution}{exo:g12:equilibrium:4}
التركيب ثابت، لكن التفاعلين المباشر والعكسي يستمران
بالسرعة نفسها: \omterm{def:g7:atoms-and-molecules:molecule}{فالجزيئات} لا تكف عن التفاعل في الاتجاهين.
\end{solution}

\begin{solution}{exo:g12:equilibrium:5}
لا، فالثابت $K$ لا يتغير (\omterm{def:g10:reaction-progress-table:system}{الحالة النهائية} هي نفسها)؛ أما المدة اللازمة لبلوغ
التوازن فتقصر.
\end{solution}

\begin{solution}{exo:g12:equilibrium:6}
$\dfrac{x^2}{(2-x)^2} = 4$، إذن $\dfrac{x}{2-x} = 2$
و$x = \qty{1.33}{mol}$: \qty{0.67}{mol} من الحمض ومثلها من \omterm{def:g11:functional-groups:alcohol}{الكحول}،
و\qty{1.33}{mol} من \omterm{def:g11:functional-groups:carboxylic-acid}{الإستر} ومثلها من الماء (الثلثان مرة أخرى).
\end{solution}

\begin{solution}{exo:g12:equilibrium:7}
$Q = (0.5 \times 0.5)/(0.5 \times 0.5) = 1 < 4$: مباشر، فيتكوّن مزيد من
\omterm{def:g11:functional-groups:carboxylic-acid}{الإستر}.
\end{solution}

\begin{solution}{exo:g12:equilibrium:8}
نحو \qty{0.52}{mol} من جهة الحمض و\qty{0.79}{mol} من جهة
\omterm{def:g11:functional-groups:carboxylic-acid}{الإستر}. وبعد نحو \qty{100}{h} يقع المنحنيان كلاهما على الخط المتقطع:
التوازن.
\end{solution}

\begin{solution}{exo:g12:equilibrium:9}
إذا تحلمأ $y$ mol: $\dfrac{y^2}{(1-y)^2} = \dfrac{1}{4}$، إذن
$\dfrac{y}{1-y} = \dfrac{1}{2}$، $y = \frac{1}{3}$: $\frac{2}{3}$ mol من
\omterm{def:g11:functional-groups:carboxylic-acid}{الإستر} ومثلها من الماء، و$\frac{1}{3}$ mol من الحمض ومثلها من \omterm{def:g11:functional-groups:alcohol}{الكحول}. وهو
\omterm{def:g6:pure-substances-and-mixtures:mixture}{الخليط} النهائي نفسه الذي تعطيه الأسترة.
\end{solution}

\begin{solution}{exo:g12:equilibrium:10}
الجسم الصلب طور نقي، لا يتغير «تركيزه»: فيُستبعد
من $Q$. وإضافة مزيد من الصلب لا تغيّر $Q$، فلا
تزيح التوازن (ما دام بعض الصلب موجودًا).
\end{solution}

\begin{solution}{exo:g12:equilibrium:11}
يحدد $K$ النسبة بين ثنائي أكسيد الكربون في الغاز وثنائي أكسيد الكربون في الماء.
فإذا فُتحت القارورة، أفلت الغاز الذي فوق الماء إلى \omterm{def:g4:air-a-mixture-of-gases:air}{الهواء}: فينخفض ثنائي أكسيد الكربون
في الغاز، ويصير $Q < K$، فيغادر مزيد من الغاز الماء، حتى لا يكاد
يبقى منه شيء.
\end{solution}

\begin{solution}{exo:g12:equilibrium:12}
عند التوازن: الحمض \omterm{def:g11:functional-groups:alcohol}{والكحول} $\frac{1}{3}$، \omterm{def:g11:functional-groups:carboxylic-acid}{والإستر} والماء
$\frac{2}{3}$. ونزع نصف الماء يُبقي $\frac{1}{3}$:
$Q = \dfrac{(2/3)(1/3)}{(1/3)(1/3)} = 2 < 4$، مباشر. ومع تفاعل $y$
إضافية: $(2/3 + y)(1/3 + y) = 4(1/3 - y)^2$، أي
$27y^2 - 33y + 2 = 0$، $y = 0.064$: فترتفع كمية \omterm{def:g11:functional-groups:carboxylic-acid}{الإستر} إلى \qty{0.73}{mol}.
\end{solution}

\begin{solution}{exo:g12:equilibrium:13}
\omterm{def:g12:equilibrium:quotient}{خارج التفاعل} العكسي يُبدَّل فيه البسط
والمقام: $Q' = 1/Q$، إذن عند التوازن $K' = 1/K$. وللحلمأة:
$1/4 = 0.25$.
\end{solution}

\begin{solution}{exo:g12:equilibrium:14}
$Q = (2 \times 2)/(1 \times 1) = 4 = K$: \omterm{def:g6:pure-substances-and-mixtures:mixture}{فالخليط} في \omterm{def:g12:equilibrium:equilibrium-state}{حالة
توازن} أصلًا؛ ولا يتغير تركيبه.
\end{solution}

\begin{solution}{exo:g12:equilibrium:15}
$n - 0.95 = 0.9025 / (4 \times 0.05) = 4.51$، إذن $n = \qty{5.5}{mol}$ من
\omterm{def:g11:functional-groups:alcohol}{الكحول} لكل \omterm{def:g10:the-mole:amount}{مول} من الحمض.
\end{solution}

\begin{solution}{pb:g12:equilibrium:1}
\textbf{1.} \ce{CH3-COOH + CH3-CH2-OH <=> CH3-COO-CH2-CH3 + H2O}.

\textbf{2.} لأنه غير تام: فهو يتوقف عند توازن يستمر فيه
الاتجاهان.

\textbf{3.} $Q = \dfrac{n_{\text{إستر}}\, n_{\text{ماء}}}{n_{\text{حمض}}\, n_{\text{كحول}}}$.

\textbf{4.} $x_{\max} = \qty{1}{mol}$؛ $x_f = \frac{2}{3}$ mol.

\textbf{5.} $\tau = 0.67$.

\textbf{6.} الحمض \omterm{def:g11:functional-groups:alcohol}{والكحول} $\frac{1}{3}$ mol لكل منهما، \omterm{def:g11:functional-groups:carboxylic-acid}{والإستر} والماء
$\frac{2}{3}$ mol لكل منهما.

\textbf{7.} $K = \dfrac{(2/3)^2}{(1/3)^2} = 4$.

\textbf{8.} تحلمأ الثلث: \omterm{def:g11:functional-groups:carboxylic-acid}{الإستر} والماء $\frac{2}{3}$، والحمض
\omterm{def:g11:functional-groups:alcohol}{والكحول} $\frac{1}{3}$: \omterm{def:g6:pure-substances-and-mixtures:mixture}{الخليط} نفسه، $Q_{eq} = 4$.

\textbf{9.} إنه التفاعل نفسه عند درجة الحرارة نفسها، و$K$
لا يتوقف إلا على درجة الحرارة.

\textbf{10.} الحمض $1 - x$؛ \omterm{def:g11:functional-groups:alcohol}{الكحول} $3 - x$؛ \omterm{def:g11:functional-groups:carboxylic-acid}{الإستر} $x$؛ الماء $x$.

\textbf{11.} $x^2 = 4(1 - x)(3 - x) = 4(3 - 4x + x^2)$، إذن
$3x^2 - 16x + 12 = 0$.

\textbf{12.} $x = (16 \pm \sqrt{256 - 144})/6 = (16 \pm 10.58)/6$: 0.903
أو 4.43؛ ولا يقع بين 0 و1 إلا $x_f = \qty{0.903}{mol}$.

\textbf{13.} \qty{90}{\percent}.

\textbf{14.} $M = \qty{88.0}{g/mol}$: $0.903 \times 88.0 = \qty{79.5}{g}$.

\textbf{15.} كل نزع للماء يخفض بسط $Q$، الذي
يصير أصغر من $K$: فتتقدم المجموعة من جديد.

\textbf{16.} يستمر التفاعل حتى يُستهلك الحمض:
\qty{100}{\percent}.

\textbf{17.} الإيثانول أرخص، وفائضه يُفصل بسهولة عن
\omterm{def:g11:functional-groups:carboxylic-acid}{الإستر} (ويُعاد تدويره).

\textbf{18.} المدة فقط: فالتفاعل الذي لا يكاد يطلق حرارة
لا يكاد يتغير ثابته مع درجة الحرارة، والتسخين لا يفعل إلا أن
يجعله يبلغ التوازن أسرع.

\textbf{19.} نحو \qty{90}{\percent} من الحمض.
\end{solution}
